\section{Purity, a forced weight-eight word, and the local types}\label{sec:forced}

Assume that \(S\leq V_{17}\) is a \(\qc{17,3,d\geq6}\)
stabilizer, so \(\dim S=14\), \(|S|=2^{14}\), and \(\dim S^\perp=20\).
The counts \(A_j\) include products of basis generators.
This section proves that \(S\) is pure, that it contains a word of
weight eight, and that the stabilizer supported on such a word has one
of five types.

\subsection{Exact weight-enumerator certificates}\label{sec:certificates}

The weight-enumerator equations count words of each weight without
specifying the full code. The MacWilliams identities express centralizer
counts in terms of stabilizer counts, and the shadow identities give
additional nonnegative counts. A certificate is an integer combination
of these necessary relations that gives a contradiction under a stated
case assumption. Its check uses integer arithmetic alone.

The ordinary weight enumerator records only total weight. A split
enumerator also records which coordinate blocks contain the nonidentity
letters. This is needed when a specified stabilizer word or subgroup
occupies one block. The formulas below are the MacWilliams and shadow
identities applied separately to each block.

Partition the physical coordinates into blocks of sizes
\(n_1,\ldots,n_s\). Let \(A^{\rm split}_{\mathbf i}\) count stabilizer words with
split weight \(\mathbf i=(i_1,\ldots,i_s)\). Put \(B=|S|=2^{n-k}\), and
\[
 \mathcal{K}_j^{\ell}(w)=\sum_{u=\max(0,j-\ell+w)}^{\min(j,w)}
          (-1)^u3^{j-u}\binom wu\binom{\ell-w}{j-u}.
\]
Character orthogonality gives the centralizer count
\[
 |\{v\in S^\perp:\operatorname{splitwt}(v)=\mathbf j\}|
 =B^{-1}\mathcal T_{\mathbf j},\quad
 \mathcal T_{\mathbf j}=\sum_{\mathbf i}
       \left(\prod_\ell \mathcal{K}_{j_\ell}^{n_\ell}(i_\ell)\right)A^{\rm split}_{\mathbf i}.
\]
The character identity behind this formula is
\[
 \mathbf 1_{v\in S^\perp}=B^{-1}\sum_{s\in S}(-1)^{\langle v,s\rangle}.
\]
The sum is \(B\) when every word of \(S\) commutes with \(v\),
and zero otherwise. Summing it over words of a fixed split weight
produces the products of Krawtchouk polynomials above. This explains why
the transformed coefficients count actual centralizer words.
On an isotropic stabilizer the parity function is linear, by
\eqref{eq:polar}. The affine set
\[
 \mathcal A=\{v:\langle v,s\rangle=\wt(s)\bmod2
                          \text{ for every }s\in S\}
\]
is a nonempty coset of \(S^\perp\), since every linear functional on
\(S\) extends to the ambient space. Character orthogonality gives its
split weight counts as \(B^{-1}\widehat{\mathcal T}_{\mathbf j}\), where
\[
 \widehat{\mathcal T}_{\mathbf j}=\sum_{\mathbf i}(-1)^{\sum i_\ell}
       \left(\prod_\ell \mathcal{K}_{j_\ell}^{n_\ell}(i_\ell)\right)A^{\rm split}_{\mathbf i}.
\]
Both transforms are nonnegative. Centralizer inclusion and distance give
\[
 \mathcal T_{\mathbf j}\geq BA^{\rm split}_{\mathbf j},\qquad
 \mathcal T_{\mathbf j}=BA^{\rm split}_{\mathbf j}\quad(1\leq\textstyle\sum j_\ell<d).
\]
These statements hold without assuming purity.

If a specified isotropic subgroup \(R_i\) lies on block \(i\), enumerate
every coset of \(R_i\) in \(R_i^\perp\). Its pattern
\(p=(p_0,\ldots,p_{n_i})\) counts words at each physical weight.
Let \(y\geq0\) count stabilizer cosets aggregated by their block
patterns and suffix weight. Then a split count is the corresponding
linear combination of products of the \(p_j\)'s.
With no prescribed subgroup a block may instead use the split counts
themselves as variables. Every actual stabilizer containing the prescribed
subgroups gives such nonnegative variables, since it is closed under their
addition. The model need not reconstruct a linear code from an arbitrary
feasible point; only necessity is used.

The row families are:
\begin{center}\small
\begin{tabularx}{\linewidth}{@{}lX@{}}\toprule
Row&Definition\\\midrule
Unit&\(A^{\rm split}_{\mathbf0}=1\).\\
Total&\(\sum A^{\rm split}_{\mathbf i}=B\).\\
Low distance&\(\mathcal T_{\mathbf j}-BA^{\rm split}_{\mathbf j}=0\), \(1\leq\sum j_\ell<d\).\\
Inclusion&\(BA^{\rm split}_{\mathbf j}-\mathcal T_{\mathbf j}\leq0\).\\
Centralizer&\(-\mathcal T_{\mathbf j}\leq0\).\\
Shadow&\(-\widehat{\mathcal T}_{\mathbf j}\leq0\).\\
Weight assumptions&\(\sum_{\sum i_\ell=w}A^{\rm split}_{\mathbf i}=A_w\), or its explicitly assumed lower bound.\\
Coset relations&A specified linear relation valid separately on every physical pattern.\\\bottomrule
\end{tabularx}
\end{center}
The checker verifies each coset relation against the regenerated patterns.
No property of a later shortened code is inserted into these reductions.
The variables may be allowed to take nonnegative real values. This
enlarges the set of possible enumerators, so infeasibility of the
resulting linear system still excludes every actual code. Integrality
is used separately in the explicitly stated parity and counting steps.

Write the system as \(Ez=b\), \(Gz\leq g\), \(z\geq0\).
A certificate consists of arbitrary integer equality multipliers \(\nu\)
and nonnegative integer inequality multipliers \(\lambda\) with
\[
 c=\nu E+\lambda G\geq0,\qquad \nu b+\lambda g<0.
\]
The condition \(c\geq0\) is checked coordinate by coordinate: since
all variables in \(z\) are nonnegative, it makes \(cz\geq0\).
Equality rows may be multiplied by either sign; inequality rows need
nonnegative multipliers so their directions are preserved.
For a feasible \(z\), one would have
\(0\leq cz\leq\nu b+\lambda g<0\), a contradiction.
The checker tests every coefficient column and the negative right-hand side
with exact integers. The multipliers may be found by numerical optimization; the certificate
is checked using integer arithmetic alone.

\subsection{Purity and the forced word}

\begin{proposition}[Exact enumerator consequences]\label{prop:17-enumerator}
The hypothetical code is pure to distance six, contains a stabilizer word
of weight eight, and satisfies
\begin{equation}\label{eq:light-budget}
 A_6+A_7\leq2.
\end{equation}
\end{proposition}
\begin{proof}
The split-coset integer certificates for stabilizer weights one through
five exclude all five possibilities without assuming purity.
For a proposed word of weight \(w\), local Clifford operations normalize
it to \(Z^w\); the system uses the coset-pattern columns of
\(\langle Z^w\rangle\) on its support, with low-distance equations
through weight five. Hence \(A_1=\cdots=A_5=0\), which implies purity
because the logical distance is at least six. With these zeros imposed,
the ordinary enumerator system plus \(A_8=0\) has an exact infeasibility
certificate with negative right-hand side \(-3,932,160\). The six
certificates are:
\begin{center}\small
\begin{tabular}{@{}lrr@{}}\toprule
Case&Columns&Negative right-hand side\\\midrule
Weight 1&17&\(-786,432\)\\
Weight 2&48&\(-5,324,800\)\\
Weight 3&45&\(-442,368\)\\
Weight 4&84&\(-81,920\)\\
Weight 5&78&\(-114,688\)\\
Pure, \(A_8=0\)&18&\(-3,932,160\)\\\bottomrule
\end{tabular}
\end{center}
``Columns'' means the number of nonnegative variables in the expanded
system.

For the last assertion, write \(\mathcal{K}_j(w)=\mathcal{K}_j^{17}(w)\)
and \(B=2^{14}\). The low-weight rows are
\(\sum_w \mathcal{K}_j(w)A_w-BA_j=0\) for \(1\leq j\leq5\), and the
shadow rows are \(-\sum_w(-1)^w\mathcal{K}_j(w)A_w\leq0\).
Combine the unit row \(A_0=1\) with multiplier \(-1,835,008\),
the total row with multiplier 116, the five low-weight rows with
\((60,25,11,3,1)\), the five pure-zero rows \(A_j=0\)
(\(1\leq j\leq5\)) with
\((-851968,-278528,-507904,-147456,-180224)\), and the shadow rows
\(j=0,\ldots,4\) with \((88,32,18,4,2)\).
The result is
\(32768(A_6+A_7)\leq65536\). Dividing by 32768 proves
\eqref{eq:light-budget}.
\end{proof}

\subsection{A nine-coordinate dimension bound}

The following additive bound makes the forced support useful. It uses a
parity property that is not visible in the first few nonnegative
MacWilliams coefficients alone.

\begin{lemma}[A nine-coordinate dimension bound]\label{lem:nine-bound}
An additive code on nine sites with minimum distance at least six has
binary dimension at most six.
\end{lemma}
\begin{proof}
A code of larger dimension contains a dimension-seven subcode. Suppose
\(|C|=128\), with nonzero weight counts \(A_6,A_7,A_8,A_9\) summing to
127. The first three nonnegative dual MacWilliams coefficients give
\[
 A_7+2A_8+3A_9\leq102,\qquad
 A_8+3A_9\geq75,\qquad
 3(A_7+A_8)-4A_9\geq17.
\]
The first two imply \(A_7\leq3A_9-48\), hence \(A_9\geq16\), and
also \(A_7+A_8\leq27\). The third then gives \(A_9\leq16\).
Equality forces
\((A_6,A_7,A_8,A_9)=(84,0,27,16)\).

On a binary additive space, weight parity is a quadratic form: its failure
to be additive is the symplectic product in \eqref{eq:polar}.
For a quadratic form \(\varepsilon\) on a \(c\)-dimensional binary space \(U\),
define its Gauss sum by
\[
 G(\varepsilon)=\sum_{v\in U}(-1)^{\varepsilon(v)}.
\]
For weight parity, this is the number of even-weight words minus the
number of odd-weight words. If the polar form has rank \(2\rho\), then
\(G(\varepsilon)\) is zero or \(\pm2^{c-\rho}\). Indeed, squaring the sum and
substituting \(y=x+z\) gives
\[
 G(\varepsilon)^2=2^c\sum_{z\in\rad(U)}(-1)^{\varepsilon(z)}.
\]
The restriction of \(\varepsilon\) to the radical is linear. The last sum is
therefore zero if that restriction is nonzero, and otherwise equals
\(|\rad(U)|=2^{c-2\rho}\). This proves the asserted possibilities.
The forced distribution would instead give Gauss sum
\(128-2\cdot16=96\). Since \(c=7\) and \(0\leq\rho\leq3\),
the possible nonzero absolute values are \(128,64,32,16\), none of
which is 96. This proves the bound.
\end{proof}

\subsection{The five local types}\label{sec:types}

Choose a weight-eight stabilizer word \(p\), put \(T=\operatorname{supp}p\),
and normalize \(p=Z^8\). Let \(R=S\cap V_T\), \(r=\dim R\).
By \eqref{eq:dimC}, the shortened centralizer has length nine and
dimension \(4+r\); by purity its minimum distance is at least six.
Lemma~\ref{lem:nine-bound} therefore gives \(1\leq r\leq2\).
The rest of the argument applies to every choice of the weight-eight
word \(p\).

If \(r=1\), then \(R=\langle Z^8\rangle\).
If \(r=2\), choose \(\zeta\in R\setminus\langle p\rangle\).
Let \(\alpha\) be the number of sites where \(\zeta\) has an \(X\)-component.
It is even by commutation. Since both \(\zeta\) and \(\zeta+p\) have weight
at least six, \(8+\alpha\geq12\), so \(\alpha\in\{4,6,8\}\).
Phase maps fixing \(Z\) remove \(Y\) on those sites. On the other
\(8-\alpha\) sites, let \(\beta\) count the \(Z\)'s of \(\zeta\).
Then \(\alpha+\beta\geq6\) and \(8-\beta\geq6\).
Replacing \(\zeta\) by \(\zeta+p\) exchanges \(\beta\) with \(8-\alpha-\beta\).
Consequently \(\alpha=4\) forces \(\beta=2\); at \(\alpha=6\) we may take
\(\beta=0\) or one; and \(\alpha=8\) forces \(\beta=0\).
Site permutations now give the following list.

\begin{center}
\begin{tabular}{@{}clcc@{}}\toprule
\(\dim R\)&Generators of \(R\)&Nonzero local weights&\(\dim C\)\\\midrule
1&\(Z^8\)&8&5\\
2&\(Z^8,\ XXXXZZII\)&6,6,8&6\\
2&\(Z^8,\ XXXXXXII\)&6,8,8&6\\
2&\(Z^8,\ XXXXXXZI\)&7,7,8&6\\
2&\(Z^8,\ X^8\)&8,8,8&6\\\bottomrule
\end{tabular}
\end{center}

This list covers every possible local stabilizer, with \(R=S\cap V_T\).
In every case \(C\) is a nine-site additive code of minimum distance at
least six, of binary dimension five for the rank-one type and six for
the rank-two types. Section~\ref{sec:classification} classifies these
codes: there are 32,865 dimension-five classes and 56 dimension-six
classes.
